\documentclass{article}
\usepackage{amsmath}
\usepackage{amssymb}

\title{Fowler's theorem for involutions}
\author{Sylvain Cappell, S. Weinberger, and M. Yan}

\begin{document}

\maketitle

Fowler, in his Ph.D. thesis, proved that if $\Gamma$ is a uniform lattice in a real semisimple group with odd torsion in $\Gamma$ then there is no compact closed manifold $M$ whose universal cover is rationally acyclic. A proof can be found in [W2]. We show that the same is true for $\Gamma$ with 2-torsion.

Without loss of generality (by considering a normal subgroup of finite index), it suffices to prove this for the special case where $\Gamma = \pi \rtimes \mathbb{Z}_{2}$ for a torsion free group $\pi$, a lattice in $G$, for which there is an involution on $M = K\backslash G/\pi$ (by isometries with the locally symmetric metric) whose fixed set $F$ is not empty. ($F$ might be disconnected; for simplicity we will write what follows just for the connected case -- there are no differences in the general case.)

Now suppose that $X^{m}$ is a manifold with fundamental group $\Gamma$, $Y$ its 2-fold cover, and suppose that the universal cover of $X$ (and therefore $Y$) are rationally acyclic. We will consider the symmetric signatures of $Y$ in the (symmetric = quadratic L-group) $L(\mathbb{R}\pi)$, where $\mathbb{R}$ is the real numbers. There is an equivalence $f:Y \to M$ which (while not degree one) gives an equivalence of symmetric signatures (because over $\mathbb{R}$, all degrees have square roots, so the symmetric signature is only sensitive to the sign of the degree of the map). Since the Novikov conjecture is true for $\pi$, the assembly map from $H_{m}(B\pi; L(\mathbb{R})) \to L_{m}(\mathbb{R}\pi)$ is injective, and this detects in the degree $m$ piece $H_{m}(B\pi; \mathbb{Z})$ the class that these manifolds represent in group homology. It follows that this map is degree one. $f_{*}[Y] = [M]$.

Now we use a cobordism argument from [W1]. We now consider the image of the fundamental class of any manifold $Z$ with fundamental group $\pi$ involution inducing this automorphism of $\pi$ and the image of $[Z]$ in $H_{m}(B\Gamma; \mathbb{Z}_{2})$. It follows from standard equivariant homotopy theory that $Z$ has an equivariant map, $g$, to $M$, and thus there is a map from its fixed set $Z^{\mathbb{Z}_2} \to F$. We claim that $g_{*}[Z] = g_{*}[Z^{\mathbb{Z}_2}]$ where we make use of the map from $\mathbb{Z}_{2} \times \pi_{1}F \to \Gamma$ (and the periodicity on the group homology of $\mathbb{Z}_{2}$ to raise the dimension from that of $F$ to $\dim M$).

This cobordism is between $Z$ and a projective space bundle over $Z^{\mathbb{Z}_2}$ -- namely the projectivized normal bundle to $Z^{\mathbb{Z}_2}$. (The fundamental class of the latter is the desired element by the Leray-Hirsch theorem.) It is explicitly $Z \times [0,1]$ and on $Z \times \{1\}$ mod out in the complement of the equivariant regular neighborhood of $Z^{\mathbb{Z}_2}$ the $\mathbb{Z}/2$ action.

Thus for $Y$, this image is 0, since the action is free. For $M$ however, this is always nonzero. The action by $\mathbb{Z}_{2}$ by isometries has fixed set which is aspherical and indeed the Borel construction for the action on $M$ shows that $\mathbb{Z}_{2}\times F \to \Gamma$ induces an injection on homology in dimension $\dim(M/\mathbb{Z}_{2})$ (and an isomorphism in higher dimensions, see [B]). Since the fundamental class of an aspherical manifold is always nontrivial in its group homology, we have a contradiction.

\section*{References}

[B] A. Borel, A seminar on transformation groups, Princeton University Press 1960

[W1] S. Weinberger, Group actions and higher signatures II, CPAM 1987

[W2] S. Weinberger, Variations on a theorem of Borel, Cambridge University Press 2022

\end{document}
